\chapter{朗兰兹纲领概观}

% ===== PDF p.003 =====
本讲我们想对朗兰兹纲领作一个粗略的介绍，并给出它算术意义的第一个例子。

回忆：域 $k$ 上连通、光滑、仿射的代数群 $G$ 称为\textbf{约化的}（reductive），若其基变换 $G_{\bar k}$ 到 $k$ 的代数闭包上不包含同构于加法群 $\mathbb{G}_a$ 的闭正规子群；等价地，$G_{\bar k}$ 的幂幺根（unipotent radical）平凡。

对每个素数 $p$，令
\[
\Z_p := \varprojlim_n \Z/p^n \Z,\qquad \Q_p := \Z_p\otimes_\Z \Q
\]
分别为 $p$ 进整数环与 $p$ 进数域。令
\[
\Af := \Bigl\{(x_p)_p\in\prod_p \Q_p \Bigm| x_p\in\Z_p \text{ 除有限多个素数 } p \text{ 外}\Bigr\}
\]
为 $\Q$ 的\textbf{有限阿代尔环\idx{阿代尔环}}（ring of finite adèles），并令
\[
\A := \Af\times\R
\]
为 $\Q$ 的\textbf{完全阿代尔\idx{阿代尔}环}。

回忆 $\A$ 是局部紧拓扑环。即：$\Af$ 上存在唯一的环拓扑使子环 $\prod_p\Z_p\subseteq\Af$ 为开；$\A$ 上则取它与 $\R$ 上通常拓扑的积拓扑。一个重要性质是：子环 $\Q\subseteq\A$（对角嵌入）上的子空间拓扑是离散的。

朗兰兹纲领的起点是 $\Q$ 上的一个约化群 $G$，最重要的例子是 $G=\GL_n$（某个 $n\ge1$）。

上述事实蕴涵：$G(\A)\cong G(\Af)\times G(\R)$ 自然地是局部紧拓扑群，且 $G(\Q)\subseteq G(\A)$ 是离散子群。具体地，取嵌入 $G\subseteq\Spec(\Q[X_1,\dots,X_m])$（某个 $m\ge1$），并取诱导嵌入 $G(\A)\subseteq\A^m$ 的子空间拓扑；这个拓扑与嵌入的选取无关。细节见 \texttt{[GH19, Theorem 2.2.1.]}。

我们将使用如下事实：

\begin{itemize}
\item 任意局部紧拓扑群 $H$ 上都存在右不变 Haar 测度，见 \texttt{[GH19, Section 3.2.]}，且它在 $\R_{>0}$ 中的正标量倍意义下唯一。
\item $G(\A)$ 上的右不变 Haar 测度同时也是左不变测度，并下降为 $G(\Q)\backslash G(\A)$ 上的 $G(\A)$-不变测度，见 \texttt{[GH19, Lemma 3.5.4.]}、\texttt{[GH19, Lemma 3.5.3]}。
\end{itemize}

现在引入朗兰兹纲领中关于 $G$ 的最终感兴趣的对象。

\noindent\textbf{定义 2.1}\phantomsection\label{def:2.1}\quad 令
\[
L^2(G(\Q)\backslash G(\A))
\]
为可测函数 $f:G(\Q)\backslash G(\A)\to\C$ 且
\[
\int_{G(\Q)\backslash G(\A)}|f|^2<\infty
\]
构成的空间，

% ===== PDF p.004 =====
其中积分关于 $G(\Q)\backslash G(\A)$ 上的 $G(\A)$-不变测度。

与通常一样，若两个可测函数在一个零测集之外相等，则在 $L^2$ 空间中视为同一个元素。

空间 $L^2(G(\Q)\backslash G(\A))$ 通过内积
\[
(f_1,f_2)\mapsto\int_{G(\Q)\backslash G(\A)}f_1(g)\overline{f_2(g)}\,\mathrm{d}g
\]
自然成为 Hilbert 空间。$G(\A)$ 经右平移作用在 $G(\Q)\backslash G(\A)$ 上，由此诱导出 $G(\A)$ 在 $L^2(G(\Q)\backslash G(\A))$ 上的（左）作用；由于所选取的测度是右不变的，这个作用是酉作用，即保持内积。

现在我们可以（以极为粗略的形式）陈述朗兰兹纲领的一个主要目标：

\begin{quote}
把 Hilbert 空间 $L^2(G(\Q)\backslash G(\A))$ 作为 $G(\A)$-表示，按算术数据加以分解。
\end{quote}

我们会略微说明「分解」是什么意思。记
\[
\widehat{G(\A)}
\]
为（Hilbert 空间上的）不可约酉 $G(\A)$-表示的同构类集。由于分析上的问题，期望存在形如
\[
L^2(G(\Q)\backslash G(\A))\ \text{“}\cong\text{”}\ \bigoplus_{[\pi]\in\widehat{G(\A)}}\pi^{\oplus m_\pi}
\]
的直和分解（其中 $m_\pi\in\N\cup\{\infty\}$ 是某个不可约酉 $G(\A)$-表示 $\pi$ 的重数）是过于天真的。

让我们考察 $G=\Gm=\GL_1$ 的例子。此时
\[
G(\Q)\backslash G(\A)\cong\Q^\times\backslash\A^\times
\]
是局部紧 abelian 群，因此可以用（抽象的）Fourier 理论来分析 $\GL_1(\A)=\A^\times$-表示
\[
L^2(\Q^\times\backslash\A^\times).
\]
下面回忆抽象 Fourier 理论的一些结果，依 \texttt{[DE14, Chapter 3]}。

\begin{itemize}
\item 设 $A$ 为任意局部紧 abelian 群（如 $\R$、$S^1\cong\R/\Z$、$\A^\times$、$\Z/n$、$\Q_p$ 等）。
\item 设 $\widehat A$ 为 $A$ 的不可约酉表示的同构类集。
\item 则 $\widehat A$ 与酉连续特征 $A\to S^1$ 的集合 $\mathrm{Hom}_{\mathrm{cts}}(A,S^1)$ 一一对应；也就是说，每个 $\pi\in\widehat A$ 都是一维的，并且（通过逐点乘法与紧开拓扑）本身又成为局部紧 abelian 群。
\item 此外 $A\cong\widehat{\widehat A}$ 作为拓扑群（同构由 $a\mapsto(\chi\mapsto\chi(a))$ 给出），并且 $A$ 离散当且仅当 $\widehat A$ 紧。\znoteplain{原文脚注：“$\Rightarrow$” 取 $\coprod_I\Z\twoheadrightarrow A$，则 $\widehat A$ 嵌入 $\prod_I S^1$。“$\Leftarrow$” 用 $A\cong\widehat{\widehat A}$ 与紧开拓扑的定义。}
\end{itemize}

% ===== PDF p.005 =====
Plancherel 定理是抽象 Fourier 理论的主要结果。对我们而言，它的意义在于给出了空间 $L^2(A)$ 作为 $A$-表示的另一种描述。

\noindent\textbf{定理 2.2}（Plancherel 定理，见 \texttt{[DE14, Chapter 3.4]}）\phantomsection\label{thm:2.2}\quad Fourier 变换
\[
\mathcal F:L^1(A)\cap L^2(A)\longrightarrow L^2(\widehat A),\qquad
f\mapsto\Bigl(\chi\mapsto\int_A f(x)\chi(x)\,\mathrm{d}\mu(x)\Bigr)
\]
延拓为 Hilbert 空间的酉同构
\[
\mathcal F:L^2(A)\xrightarrow{\ \sim\ }L^2(\widehat A).
\]
当我们让 $A$ 通过
\[
a\cdot g(\chi):=\chi(a)g(\chi),\qquad a\in A,\ g\in L^2(\widehat A),\ \chi\in\widehat A
\]
作用在 $L^2(\widehat A)$ 上时，Fourier 变换 $\mathcal F$ 是 $A$-等变同构。因此作为 $A$-表示有
\[
L^2(A)\cong L^2(\widehat A).
\]

让我们讨论两个例子：$A=S^1$ 与 $A=\R$。

\begin{itemize}
\item 若 $A=S^1$，则
\[
\Z\cong\widehat{A},\qquad n\mapsto\chi_n,\quad \chi_n(z)=z^n,\ z\in S^1.
\]
于是
\[
L^2(S^1)\cong\widehat{\bigoplus_{n\in\Z}}\C\chi_n
\]
是 $S^1$-等变子空间 $\C\chi_n\subseteq L^2(S^1)$ 的 Hilbert 空间直和。具体地说：每个 $L^2$ 函数 $f:S^1\to\C$ 都可以唯一地写成
\[
f=\sum_{n\in\Z}a_n\chi_n,\qquad a_n\in\C,\quad \sum_{n\in\Z}|a_n|^2<\infty.
\]
这样我们就找到了酉 $S^1$-表示 $L^2(S^1)$ 分解为不可约子空间的一个漂亮分解，分析上的问题并不大。不过，我们不得不考虑 Hilbert 空间直和 $\widehat{\bigoplus}_{n\in\Z}\C\chi_n$，而不是（代数的）表示直和。
\item 若 $A=\R$，则 $\widehat{\R}\cong\R$，同构由
\[
\R\to\widehat{\R},\qquad x\mapsto\chi_x,
\]
给出，其中 $\chi_x:\R\to S^1$，$y\mapsto e^{2\pi ixy}$。同构
\[
\mathcal F:L^2(\R)\xrightarrow{\ \sim\ }L^2(\R)
\]
因而蕴涵：每个 $f\in L^2(\R)$ 都可以写成
\[
f(x)=\int_\R g(y)\chi_y(x)\,\mathrm{d}y
\]

% ===== PDF p.006 =====
的形式，其中 $g\in L^2(\R)$ 唯一。因此
\[
L^2(\R)\cong\int_\R\chi_y\,\mathrm{d}y
\]
是表示的\textbf{Hilbert 积分}（见 \texttt{[GH19, Section 3.7]}），这个概念不像 Hilbert 空间直和那样直接。注意我们无法做得更好：对任意 $x\in\R$ 都有 $\chi_x\notin L^2(\R)$，而且 $L^2(\R)$ 不含 $\R$ 的任何不可约子表示，见 \texttt{[GH19, Section 3.7]}。\znoteplain{原文脚注：也没有任何不可约商；\texttt{[GH19, Lemma 3.7.1]} 中关于商的陈述，我认为是错的。}
\end{itemize}

让我们回到 $G=\Gm$，即 $A=\A^\times$ 的情形。

\begin{itemize}
\item 把每个 $n\in\Q^\times$ 分解为 $n=\pm p_1\cdots p_k$（$p_i$ 为素数），可推出
\[
\Q^\times\backslash\A^\times\cong\prod_p\Z_p^\times\times\R_{>0}.
\]
\item 我们暂时忽略因子 $\R_{>0}$，考察
\[
L^2\Bigl(\prod_p\Z_p^\times\Bigr)\cong\widehat{\bigoplus_{\chi}}\C\chi,
\]
其中 $\chi:\prod_p\Z_p^\times\to S^1$ 遍历 $\prod_p\Z_p^\times$ 的全部连续特征。注意每个特征 $\chi:\prod_p\Z_p^\times\cong\varprojlim_m(\Z/m)^\times\to S^1$ 都通过某个商 $(\Z/m)^\times$ 分解。
\item 这些特征 $\chi$ 是「算术」数据，即
\[
\Gal(\Q(\mu_\infty)/\Q)\cong\prod_p\Z_p^\times,
\]
其中 $\Q(\mu_\infty)=\bigcup_{n\in\N}\Q(e^{2\pi i/n})$ 是 $\Q$ 的分圆扩张。定义
\[
\alpha:\Gal(\overline\Q/\Q)\twoheadrightarrow\Gal(\Q(\mu_\infty)/\Q)\cong\prod_{p\ \text{素数}}\Z_p^\times
\]
为自然投影。
\item 因此对每个 $\chi:\prod_p\Z_p^\times\to S^1$，复合 $\chi\circ\alpha$ 都是 $\Gal(\overline\Q/\Q)$ 的一个（连续）一维 Galois 表示。
\item 由 Kronecker--Weber 定理，每个一维 Galois 表示
\[
\sigma:\Gal(\overline\Q/\Q)\to\C^\times
\]
都具有 $\sigma=\chi\circ\alpha$ 的形式，其中 $\chi$ 如上。
\end{itemize}

\noindent\textbf{定理 2.3}（Kronecker--Weber）\phantomsection\label{thm:2.3}\quad 域 $\Q(\mu_\infty)$ 是 $\Q$ 的极大 abelian 扩张；也就是说，每个 Galois 群 $\Gal(F/\Q)$ 为 abelian 的有限 Galois 扩张 $F/\Q$ 都包含在 $\Q(\mu_\infty)$ 中。

于是对 $G=\Gm$ 我们得到如下分解：
\[
L^2(\R_{>0}\Q^\times\backslash\A^\times)\cong\widehat{\bigoplus_{\sigma:\Gal(\overline\Q/\Q)\to\C^\times}}\C\chi_\sigma,
\]
% ===== PDF p.007 =====
其中 $\chi_\sigma$ 由 $\chi_\sigma\circ\alpha=\sigma$ 刻画。这个（简单的）例子是朗兰兹纲领所追求目标的原型。注意模掉 $\R_{>0}$ 是方便而且无害的。全空间
\[
L^2(\Q^\times\backslash\A^\times)
\]
是表示 $L^2(\R_{>0}\Q^\times\backslash\A^\times)$ 在 $\R_{>0}$ 的酉特征空间 $\widehat{\R}\cong\R$ 上的 Hilbert 积分。

现在转向一般的 $G$，讨论空间
\[
G(\Q)\backslash G(\A)
\]
的几何。这将比 $G=\GL_1$ 的情形复杂得多。

存在中心子群 $A_G\subseteq G(\R)$，$A_G\cong\R_{>0}^r$，使得
\[
[G]:=A_GG(\Q)\backslash G(\A)
\]
具有有限体积，见 \texttt{[GH19, Theorem 2.6.2.]}。具体地，$A_G$ 可以取作 $G$ 的中心 $Z(G)$ 的极大分裂子环面的连通分支；例如 $G=\GL_n$ 时 $A_{\GL_n}\subseteq\GL_n(\R)$ 是元素在 $\R_{>0}$ 中的标量矩阵构成的群，而 $A_{\SL_n}$ 平凡。

同样，全空间 $L^2(G(\Q)\backslash G(\A))$ 是表示 $L^2([G])$ 在 $\widehat{A_G}$ 上的 Hilbert 积分，因此用 $L^2([G])$ 代替 $L^2(G(\Q)\backslash G(\A))$ 可以省去一些分析上的麻烦。\znoteplain{原文脚注：精确地说，$G(\A)$ 上每个右不变 Haar 测度都下降为 $[G]$ 上的测度，见 \texttt{[GH19, Section 2.6]}。}

虽然我们只关心 $G(\A)$ 上任意的右不变测度\znoteplain{原文脚注：因为所得的 $L^2([G])$ 空间彼此同构。}，但请允许我们提到：$G(\A)$ 上存在一个典范测度，即\textbf{玉河测度}（Tamagawa measure），对这个测度而言，数
\[
\tau(G):=\mathrm{vol}([G])
\]
称为 $G$ 的\textbf{玉河数\idx{玉河数}}（Tamagawa number）。玉河数的知识在算术上很有意义，参见 \texttt{[Col11, Appendix B]}、\texttt{[Lur, Lecture 1]}。例如 $\tau(\SL_n)=1$，这蕴涵
\[
\mathrm{vol}(\SL_n(\R)/\SL_n(\Z))=\zeta(2)\zeta(3)\cdots\zeta(n),
\]
其中 $\zeta(s)$ 是 Riemann $\zeta$ 函数。

拓扑空间 $G(\Q)\backslash G(\A)$ 与 $[G]$ 都是某个实流形的 profinite 覆盖。例如 $G=\GL_n$ 时

% ===== PDF p.008 =====
$\GL_n(\Z)\backslash\bigl(\prod_p\GL_n(\Z_p)\times\GL_n(\R)\bigr)$（$\GL_n(\Z)$ 对角地作用）覆盖实流形 $\GL_n(\Z)\backslash\GL_n(\R)$；类似地 $[G]$ 覆盖 $\GL_n(\Z)\R_{>0}\backslash\GL_n(\R)$。见 \texttt{[GH19, Section 2.6]}、\texttt{[Del73, (0.1.4.1)]}。

更精确地说，对每个紧开子群 $K\subseteq G(\Af)$（也称为一个「level 子群」），商
\[
[G]/K
\]
都是实流形。

类似地，设 $K_\infty\subseteq G(\R)$ 为 $G(\R)$ 的极大紧子群（连通与否不论），例如 $G=\GL_n$ 时可取 $K_\infty\cong\SO_n(\R)$。由于 $K_\infty$ 与 $A_G\subseteq G(\R)$ 平凡相交，空间 $[G]$ 是 $[G]/K_\infty$ 上的主 $K_\infty$-丛（或称 $K_\infty$-挠子）。

下一讲我们将更详细地讨论 $G=\GL_2$ 的情形。现在只提一点：若 $K_\infty\cong\SO_2(\R)$，则由 Möbius 变换给出同构
\[
A_G\backslash\GL_2(\R)/K_\infty\xrightarrow{\ \sim\ }\mathbb{H}^\pm,\qquad
g=\begin{pmatrix}a&b\\c&d\end{pmatrix}\mapsto g(i):=\frac{a\cdot i+b}{c\cdot i+d},
\]
其中
\[
\mathbb{H}^\pm:=\{z\in\C\mid \operatorname{Im}(z)\neq0\}
\]
是上半平面与下半平面之并。特别地，
\[
[\GL_2]/K_\infty\cong\GL_2(\Z)\backslash\Bigl(\GL_2\Bigl(\prod_p\Z_p\Bigr)\times\mathbb{H}^\pm\Bigr).
\]

设 $K\subseteq G(\Af)$ 为紧开子群，则集合
\[
G(\Q)\backslash G(\Af)/K=\bigsqcup_{i=1}^m G(\Q)g_iK
\]
有限（见 \texttt{[GH19, Theorem 2.6.1]}），并且
\[
[G]/K_\infty K\cong\bigsqcup_{i=1}^m\Gamma_i\backslash X,
\]
其中
\[
X:=A_G\backslash G(\R)/K_\infty
\]
是实流形，而
\[
\Gamma_i:=G(\Q)\cap G(\R)g_iKg_i^{-1}\subseteq G(\R)
\]
是离散子群（同余子群的一个例子，见 \texttt{[GH19, Section 2.6]}）。当 $K$ 充分小时，$\Gamma_i$ 在 $X$ 上自由且真不连续地作用，见 \texttt{[GH19, Definition 15.2]}。此时态射 $X\to\Gamma_i\backslash X$ 是以 $\Gamma_i$ 为覆盖群的（拓扑意义下的）覆盖。我们最关心的情形是 $G=\GL_2$，此时一切都变得更具体：\[\Gamma_i\backslash X\cong\Gamma\backslash\mathbb{H}^\pm,\]

其中 $\Gamma\subseteq\GL_2(\Z)$ 是包含
\[
\Gamma(m):=\ker(\GL_2(\Z)\to\GL_2(\Z/m))
\]
（某个 $m\ge0$）的离散子群。当 $m\ge3$ 时，$\Gamma(m)$ 在 $\mathbb{H}^\pm$ 上的作用自由。这个例子将在下一讲提供与模形式的联系。\znote{原文此处作 “For K is sufficiently small”，语法有误，译文按 “When $K$ is sufficiently small” 理解。}

注意对一般的 $G$，通常 $\dim(\Gamma_i\backslash X)>0$。因此 $[G]$ 的几何、以及空间 $L^2([G])$，都比 $G=\Gm$ 的情形复杂得多。

实流形 $\Gamma_i\backslash X$ 本身是很有趣的实流形，因为它们具有算术意义。我们将会看到，当 $G=\GL_2$ 时这些 Riemann 曲面是拟射影的，并典范地定义在数域上。

若空间 $[G]$ 紧（这当且仅当任何 $\Gm\subseteq G$ 都落在中心内，见 \texttt{[GH19, Theorem 2.6.2]}），则 $L^2([G])$ 的分析更简单。一个具体的例子是 $G$ 取 $\Q$ 上非分裂四元数代数的单位群，即具有表现

% ===== PDF p.009 =====
\[
\langle x,y\mid x^2=a,\ y^2=b,\ xy=-yx\rangle
\]
（适当选取 $a,b\in\Q^\times$）的 $\Q$-代数的单位群。注意当 $ab<0$ 时 $G(\R)\cong\GL_2(\R)$，因此 $\Gamma_i\backslash X$ 仍然是上半平面的商。四元数代数的更多细节见 \texttt{[GS17, Chapter 1]}。

$\Q$ 上的约化群大量存在。产生大量例子的一种特别容易的方式是通过 Weil 限制。为此设 $F/\Q$ 为有限扩张，$H$ 在 $F$ 上约化。则函子 $R\mapsto H(R\otimes_\Q F)$（定义在 $\Q$-代数范畴上）可由 $\Q$ 上的约化群
\[
G:=\Res_{F/\Q}(H)
\]
表出，见 \texttt{[CGP15, Appendix A]} 或 \texttt{[GH19, Section 1.4]}。对这个约化群 $G$ 有
\[
G(\Q)\backslash G(\A)=H(F)\backslash H(\A_F),
\]
其中 $\A_F$ 是 $F$ 的阿代尔环。这蕴涵：就朗兰兹纲领而言，考虑 $\Q$ 上所有约化群等价于考虑所有数域上所有约化群。

当 $H=\mathbb{G}_{m,F}$ 时，$L^2([G])$ 的期望分解（大致上）等价于 $F$ 的类域论。顺带一提：$[G]$ 体积的有限性包含了 $F$ 的类数有限性与 $F$ 上的 Dirichlet 单位定理。

另一个例子来源：设 $F$ 为虚二次域，$H=\GL_{2,F}$。此时 $G(\R)\cong\GL_2(\C)$，可取 $K_\infty\cong U(2)$（酉群），而
\[
X=A_G\backslash G(\R)/K_\infty
\]
是三维双曲空间 $\mathbb{H}^3$。此时的商 $\Gamma_i\backslash X$ 称为\textbf{算术双曲 3-流形}，见 \texttt{[Thu82, Section 4 \& 5]}。更精确地说，双曲 3-空间 $\mathbb{H}^3$ 可以模型在四元数
\[
\{q=x+y\cdot i+z\cdot j\mid x,y\in\R,\ z>0\}
\]
上，而 $\begin{pmatrix}a&b\\c&d\end{pmatrix}\in\GL_2(\C)$ 通过 $q\mapsto(aq+b)(cq+d)^{-1}$ 作用在 $\mathbb{H}^3$ 上，其中 $\C=\{x+y\cdot i\mid x,y\in\R\}$。

讨论完一些例子之后，我们现在可以给出朗兰兹纲领中居于中心地位的自守表示\idx{自守表示}（在 $L^2$ 意义下）的定义。

\noindent\textbf{定义 2.4}（\texttt{[GH19, Definition 3.3]}）\phantomsection\label{def:2.4}\quad 设 $G$ 为 $\Q$ 上的约化群。一个\textbf{自守表示}（在 $L^2$ 意义下）是不可约酉 $G(\A)$-表示 $\pi$，它同构于 $L^2([G])$ 的某个子商（subquotient）。

以后我们还会讨论另一种更代数的自守表示概念。在上述定义中，$\pi$ 必然在 $A_G$ 上平凡。这个假设通常是无害的：由不可约酉表示的 Schur 引理，对任意不可约酉 $G(\A)$-表示 $\pi$ 都能找到特征 $\chi:G(\R)\to\C^\times$，使得 $\chi\otimes\pi$ 在 $A_G$ 上平凡，见 \texttt{[GH19, Section 6.5]}。例如 $G=\GL_n$ 时总可以取 $\chi$ 为形如 $g\mapsto|\det(g)|^s$（某个 $s\in i\R$）的特征。

空间 $L^2([G])$ 分解为「离散」部分与「连续」部分，见 \texttt{[GH19, Section 10.4]}。从分析的角度看，$[G]$ 紧的情形更容易，因为此时连续部分消失。

现在我们更精确地解释「离散」部分与「连续」部分的含义。回忆 $\widehat{G(\A)}$ 记不可约

% ===== PDF p.010 =====
酉 $G(\A)$-表示的同构类集。类比于局部紧 abelian 群 $A$ 上 $\widehat A$ 的紧开拓扑，集合 $\widehat G$ 有一个自然的拓扑，即 \textbf{Fell 拓扑}，见 \texttt{[GH19, Section 3.8]}。粗略地说，Fell 拓扑可以这样描述：取 $\pi\in\widehat{G(\A)}$、$x\in\pi$ 以及连续 $\C$-线性同态 $\phi:\pi\to\C$，则
\[
f_{x,\phi}:G\to\C,\qquad g\mapsto\phi(gx)
\]
称为 $\pi$ 的一个\textbf{矩阵系数}。粗略地说，$\pi,\pi'\in\widehat{G(\A)}$ 在 Fell 拓扑中彼此接近，若它们的矩阵系数在紧开拓扑下彼此接近。

下面的定理对 $G(\A)$ 的酉表示更一般地成立。我们不需要定理 2.5 中积分的精确含义，只需指出它的出现是因为 $G(\A)$ 的某些表示「贡献」于 $L^2([G])$ 却不作为子表示出现（类似于 $\R$ 的 Fourier 理论中的情形）。

\noindent\textbf{定理 2.5}（\texttt{[GH19, Theorem 3.9.4]}）\phantomsection\label{thm:2.5}\quad 存在可测的重数函数 $m:\widehat G\to\{1,2,\dots,\infty\}$ 与 $\widehat G$ 上的测度 $\mu$，使得
\[
L^2([G])\cong\int_{\widehat G}\Bigl(\widehat{\bigoplus^{m(\pi)}}\pi\Bigr)\,\mathrm{d}\mu(\pi).
\]
$m$ 与 $\mu$ 在「零测集上作改变」的意义下唯一。

借助这个定理，我们现在可以粗略解释 $L^2([G])$ 的「离散」部分与「连续」部分：满足 $\mu(\pi)>0$ 的点 $\pi\in\widehat{G(\A)}$ 对 $\widehat{G(\A)}$ 上的测度 $\mu$ 而言是「离散」的，它们以重数 $m(\pi)$ 作为子表示出现在 $L^2([G])$ 中。这引出分解
\[
L^2([G])\cong L^2_{\mathrm{disc}}([G])\oplus L^2_{\mathrm{cont}}([G]),
\]
其中 $L^2_{\mathrm{disc}}([G])$ 是 $G(\A)$ 的全部不可约酉子表示的 Hilbert 和，$L^2_{\mathrm{cont}}([G])$ 是它的正交补，见 \texttt{[GH19, Section 9.1]}。

Langlands 在一项开创性的工作中用 $G$ 的（真）抛物子群\idx{抛物子群}描述了连续部分。回忆：按定义，闭连通子群 $P\subseteq G$（在 $\Q$ 上）称为\textbf{抛物子群}，若 $G/P$ 是 $\Spec(\Q)$ 上的 proper 概形；而 $P$ 对其单根取商所得为其 \textbf{Levi 商} $M$，见 \texttt{[GH19, Section 1.9]}。例如对 $\GL_n$，取分解 $n=n_1+n_2+\dots+n_k$，令 $P$ 为 $\GL_n$ 中块大小为 $n_1,\dots,n_k$ 的块上三角矩阵构成的子群；其 Levi 商为块对角矩阵构成的群。$\GL_n$ 的所有抛物子群在共轭下都呈这种形状。

Langlands 随后证明：连续部分 $L^2_{\mathrm{cont}}([G])$ 可以通过对 $P\subseteq G$（$P\neq G$）的所有 Levi 商的离散部分 $L^2_{\mathrm{disc}}([M])$ 作「诱导」来描述，见 \texttt{[GH19, Section 10.4]}。这就是他的「Eisenstein 级数\idx{Eisenstein 级数}」理论，也是朗兰兹纲领的一个起点。

对本课程而言，$L^2_{\mathrm{disc}}([G])$（实际上是其中恰当的「尖点」子空间）将比整个空间 $L^2([G])$ 更重要。顺带一提：$[G]$ 紧当且仅当 $G$ 没有（定义在 $\Q$ 上的）真抛物子群。这与前面所说的「此时 $L^2([G])=L^2_{\mathrm{disc}}([G])$」是一致的。
